\scanpage{7}
\uline{Unités} On définit
\[
\Delta_X\in\operatorname{Hom}_{\unsure{\mathcal V_c^A|\mathcal V^A}}(X,X)=A(X\times X)
\]
en considérant
\[
\delta_X:X-X\times X\qquad(\mathrm{diag})
\]
et définissant
\[
\Delta_X=\deleted{\unsure{\delta_{X*}A(\delta_X)}}A(\delta_X)_*(1_{A(X)})
\]
Prouvons que $\Delta_X$ est une unité \unsure{bilatère} pour la
composition : on a
\diagram{\begin{tikzcd}
X \arrow[r,no head,"\Delta_X"] & X \arrow[r,no head,"u"] & Y
\end{tikzcd}}
\[
u\circ\Delta_X=p_{31*}(p_{32}^*(u)\,p_{21}^*(\Delta_X))
\]
\diagram{\begin{tikzcd}[column sep=small]
&X\times X\times Y\arrow[dl,"p_{21}"']\arrow[d,no head]\arrow[dr,"p_{32}"]&\\
X\times X&X\times Y&X\times Y
\end{tikzcd}}
\[
\begin{aligned}
p_{21}^*(\Delta_X)&=\unsure{p_{21}^*|p'_{21}{}^*}(\delta_{X*}(1))
\overset{\text{\inserted{\fbox{échange}}}}=q_*(p^*(1))=q_*(1)\\
p_{32}^*(u)\,p_{21}^*(\Delta_X)&=p_{32}^*(u)\,q_*(1)=q_*(q^*(p_{32}^*(u)))
\end{aligned}
\]
\margin{form. de proj\quad Attention \uline{signe!}}
\diagram{\begin{tikzcd}[column sep=large,row sep=large]
X\times Y\arrow[r,"\delta_X\times\operatorname{id}_Y=q"]\arrow[d,"p"']&X\times X\times Y\arrow[d,no head,"p_{21}"]\\
X\arrow[r,no head,"\delta_X"]&X\times X
\end{tikzcd}}
\margin{\unsure{donc}\quad\fbox{\parbox{.35\linewidth}{signe $(-1)^{uX}$ ; pas de signe pour calculer \deleted{\unsure{que}} $\Delta_X\circ v=v$\quad $v:Y-X$}}}
\[
\deleted{\unsure{a}}\quad=q_*(u)\qquad\text{car }p_{32}q=\operatorname{id}_{X\times Y}
\]
donc
\[
u\circ\Delta_X=p_{31*}q_*(u)=u
\]
car \deleted{\unsure{q}} $p_{31}q=\operatorname{id}$

On procède de façon analogue pour $\Delta_X\circ u=u$.

\uline{Donc on trouve bien une catégorie $\mathcal V^A$}

[\uline{NB} Pour définir la composition on n'a eu besoin que
des $p_*,p^*$ pour projections $X\times Y\unsure{\rightleftarrows|\longrightarrow}X$ ; mais pour
\deleted{définir} \unsure{construire} les unités on a utilisé aussi $q_*,q^*$ pour
les imm. diagonales et leurs produits par un $Y$\ldots]
\mspage{6}
